\capitulo{6}{Trabajos relacionados}

Este apartado presenta un resumen comentado de los trabajos y proyectos más relevantes en el campo de la programación por bloques y la robótica educativa. El objetivo es contextualizar el prototipo desarrollado en este Trabajo Fin de Grado (TFG) dentro del panorama actual, identificando las tendencias, fortalezas y limitaciones de las soluciones existentes. Este análisis comparativo ha servido para fundamentar las decisiones de diseño e implementación del proyecto, buscando aportar valor y abordar necesidades específicas.

\section{Plataformas de programación por bloques}
En el ámbito de la programación visual, existen plataformas consolidadas que han servido como inspiración y punto de partida teórico para este proyecto.
\begin{itemize}
	\item \textbf{Scratch \cite{Paredes2023}:} Desarrollado por el MIT Media Lab, es la referencia mundial para la enseñanza de la programación visual a niños y jóvenes. Su principal fortaleza reside en su interfaz intuitiva, el modelo de "drag and drop" de bloques, su categorización de colores, y su enfoque en la creación de narrativas y animaciones. Si bien es extremadamente accesible, su ecosistema es relativamente cerrado y no está diseñado para integraciones directas con entornos 3D o hardware robótico específico fuera de su plataforma.
	\item \textbf{Blockly \cite{Munoz2015}:} Librería de Google de código abierto que proporciona el framework base para construir editores de programación por bloques. Destaca por su alta personalización y capacidad para generar código textual en múltiples lenguajes. Es la base técnica subyacente para numerosas soluciones, incluyendo extensiones hacia el aprendizaje de redes neuronales \cite{Wang2021_BlockDL} o el control de dispositivos IoT \cite{Ionescu2020_BlockIoT}. Aunque potente, carece de la capa de experiencia de usuario y del ecosistema lúdico inherente a Scratch.
	\item \textbf{UBlockly \cite{UBlockly}:} Como implementación de Blockly en C\# para Unity, UBlockly ha sido la \textbf{principal referencia técnica para la fase inicial de este proyecto}. Su estructura modular y su capacidad para integrarse con el motor Unity (gestionando modelos de bloques y ejecución mediante corrutinas) la convirtieron en un excelente punto de partida para evaluar la viabilidad de un sistema de programación visual. Su estructura modular y su capacidad para integrarse con el motor Unity (gestionando modelos de bloques, ejecución mediante corrutinas, y serialización) la convierten en una excelente plataforma para el desarrollo de un editor de bloques funcional y extensible.  Aunque el prototipo final evolucionó hacia un sistema híbrido, el análisis de su arquitectura (y la refactorización a MVC que se realizó sobre ella) fue fundamental para definir los componentes clave del sistema final.
\end{itemize}

\section{Simuladores y entornos de robótica educativa}
Para abordar la limitación de acceso a hardware físico, han surgido diversos simuladores que permiten a los estudiantes programar robots virtuales.
\begin{itemize}
	\item \textbf{Simuladores genéricos basados en Unity:} Existen esfuerzos para desarrollar simuladores 3D de robótica en Unity para la enseñanza de programación o inteligencia artificial \cite{Simulador3DUnity}. Estos proyectos, al igual que el presente TFG, aprovechan las potentes capacidades de renderizado 3D y el entorno de desarrollo versátil de Unity. Suelen enfocarse en una demostración de la viabilidad técnica o en contextos muy específicos (como la robótica industrial), pero a menudo carecen de una interfaz de programación por bloques totalmente integrada o de un fuerte enfoque en la educación preuniversitaria.
	\item \textbf{Plataformas web con simulación 3D:}
	\begin{itemize}
		\item \textbf{Delightex (CoSpaces Edu):} Permite crear entornos 3D y aplicaciones de Realidad Virtual (VR)/Aumentada (AR) programables con bloques o JavaScript \cite{CoSpaces2020_Edu}. Es accesible vía navegador y ofrece un entorno visual atractivo para proyectos creativos, aunque la interacción con robots simulados es a menudo menos detallada que en motores dedicados.
		\item \textbf{VEXcode VR y CoderZ:} Ofrecen entornos 3D para programar robots educativos específicos (VEX, GoPiGo) mediante bloques o lenguajes textuales. Son excelentes para aprender el currículo de robótica asociado a esos kits, pero son plataformas con licencias o enfoque más específicos que restringen la flexibilidad o la personalización de robots y escenarios, y a menudo no están optimizadas para la creación de contenido 3D personalizado de forma sencilla \cite{roboticaeducativa}.
	\end{itemize}
	\item \textbf{Simuladores para hardware específico:}
	\begin{itemize}
		\item \textbf{Robot Virtual Worlds :} Ofrece simuladores de alta fidelidad para robots LEGO EV3 o VEX. Proporcionan un alto grado de realismo físico y la capacidad de probar programas diseñados para robots específicos, siendo valiosos para entrenamientos detallados \cite{RobotVirtualWorlds}. No obstante, suelen requerir mayores recursos computacionales y están fuertemente ligados a sus marcas, limitando su adaptabilidad.
		\item \textbf{Simuladores de robótica educativa de makeblock:} Existen simuladores como Makeblock mBlock que integran un editor de bloques con entornos virtuales 2D/3D para los robots de su marca, promoviendo el pensamiento computacional \cite{roboticaeducativa_makeblock}. Al igual que otros, están centrados en el ecosistema de la marca.
	\end{itemize}
	
\end{itemize}

\section{Herramientas para el desarrollo de interfaces de bloques}
Existen recursos y frameworks específicos diseñados para ayudar en la creación de interfaces de programación por bloques.
\begin{itemize}
	\item \textbf{Blocks Engine 2 \cite{Meadowgames}} Es un paquete de assets para Unity diseñado para facilitar el desarrollo de juegos y aplicaciones con interfaces de bloques. Este motor abstrae gran parte de la complejidad de la lógica de drag and drop, conexiones y validación de bloques en Unity. Representa una alternativa al enfoque de adaptar UBlockly. Si bien ofrece una solución más rápida lista para usar para la UI de bloques, mi decisión  inicial de adaptar UBlockly y la posterior evolución a un sistema textual propio se fundamentó en la necesidad de un control total sobre el motor de ejecución y la lógica de interpretación de comandos. El objetivo del TFG era precisamente profundizar en el diseño e implementación de un intérprete de lógica de programación (ya fuera visual o textual), no solo utilizar una herramienta "plug-and-play"
\end{itemize}

\section{Síntesis de la contribución del proyecto}
Tras el análisis de las herramientas existentes, la contribución fundamental de este Trabajo Fin de Grado se centra en los siguientes puntos:

\begin{itemize}
	\item \textbf{Un puente pedagógico entre lo visual y lo textual:} A diferencia de las plataformas que son o 100\% visuales (Scratch) o que se centran en la programación textual tradicional, este proyecto propone una solución híbrida. El usuario construye su programa mediante la adición secuencial de bloques de texto, manteniendo la claridad estructural de la programación por bloques, pero al mismo tiempo practica la sintaxis de un lenguaje de comandos. Esto lo posiciona como una herramienta ideal para la transición del aprendizaje, un paso intermedio que facilita el salto de Scratch a lenguajes como Python.
	
	\item \textbf{Diseño e implementación de un lenguaje de dominio específico (DSL) y su Intérprete:} El núcleo de la aportación técnica no reside en la adaptación de una librería existente, sino en el diseño e implementación desde cero de un pequeño sistema de compilación:
	\begin{itemize}
		\item Un DSL con una gramática simple y orientada a la robótica.
		\item Un Parser que actúa como analizador léxico y sintáctico.
		\item Un motor de ejecución que interpreta la estructura de datos generada por el parser.
	\end{itemize}
	Este ciclo completo de código fuente -> análisis -> ejecución es una demostración práctica de conceptos fundamentales de la ingeniería informática.
	
	\item \textbf{Integración completa en un entorno de simulación 3D accesible:} El proyecto integra esta lógica de programación en un entorno 3D interactivo construido en Unity, con un enfoque en la accesibilidad para plataformas móviles (tablets) y sin la dependencia de hardware costoso. Se logra así un laboratorio virtual completo y autónomo, alineado con los Objetivos de Desarrollo Sostenible de democratizar la educación tecnológica.
	
\end{itemize}
En síntesis, este TFG no se limita a replicar una solución existente, sino que aporta un prototipo con un enfoque pedagógico y técnico original, abordando el nicho de la transición a la programación textual mediante un entorno de desafíos gamificado y accesible.
